% Standalone section on the order of W(n); uses the macros of version-5.tex
% (\Hs, \B, \degpm). Compile by \input-ing after the preamble of version-5.tex.
\section{On the order of the worst-case head complexity}
\label{sec:wn-order}

Write $N=2^n$ and $W(n)=\max_f\Hs(f)$. Theorems~\ref{thm:worst} and
\ref{thm:asymptotic-spanning} give
$(\tfrac12+o(1))N/n^2\le W(n)\le(1+o(1))N/n$. This section records what
is known about the order of $W(n)$ after an audit of those proofs, two
new rigidity theorems, and an analysis of the mechanisms that a
$2^n/n^2$ construction would have to use. The question whether
$W(n)=O(N/n^2)$ remains open; the section identifies precisely where
every known mechanism stops.

\subsection{Normalized scores}
\label{sec:wn-normalized}
Every $H$-head score can be written as
\begin{equation}
\label{eq:wn-normal}
 S(x)=\sigma+\sum_{h\in P}\frac{u_h\cdot x}{1+b_h\cdot x}
          +\sum_{h\in M}\frac{v_h\cdot(\mathbf1-x)}{1+d_h\cdot(\mathbf1-x)},
 \qquad b_h,d_h\in\mathbb R^n_{\ge0},
\end{equation}
with $\sigma\in\mathbb R$ and $u_h,v_h\in\mathbb R^n$ arbitrary.
Indeed, dividing an atom with positive slopes by $B_h(0)>0$ and writing
$A_h/B_h=A_h(0)+(u_h-A_h(0)b_h)\cdot x/B_h$ absorbs all numerator
constants into $\sigma$; an atom with negative slopes becomes, in the
complemented variables $y=\mathbf1-x$, an atom with positive slopes;
constant-slope atoms are affine functions and are absorbed into one
extra head by the perturbation in the proof of Theorem~\ref{thm:nf}.
For fixed denominators the realizable Boolean functions are exactly the
halfspace dichotomies of the $N$ points
$\phi_B(x)=(1,(x_i/B_h(x))_{h,i})\in\mathbb R^{1+nH}$
(Proposition~\ref{prop:uniform-dual-certificate}); in particular
Cover's count bounds the number of functions realizable with any one
fixed list by $2\sum_{i\le nH}\binom{N-1}{i}$.

\begin{proposition}[Where the information must sit]
\label{prop:wn-information}
Let $H=cN/n^2$ with $c$ fixed. For any fixed denominator list, at most
$2\sum_{i\le nH}\binom{N-1}i\le 2^{(N/n)\log_2(ecn)+O(1)}$ Boolean
functions are realizable. Consequently a representation of all Boolean
functions with $H=cN/n^2$ heads must use denominators that depend on
the target, and the $nH$ denominator parameters must carry $\Theta(n)$
bits each, while the numerator parameters carry $O(\log n)$ bits each.
\end{proposition}
\begin{proof}
The first claim is Cover's bound for a linear family of dimension
$nH+1$ on $N$ points, using $\binom{N-1}{i}\le(eN/i)^i$. For the second,
the number of realizable functions is at most the number of distinct
fixed-list pattern sets times the maximum size of one such set; the
latter is $2^{(N/n)\log_2(ecn)+O(1)}=2^{o(N)}$, so $2^{N}$ functions
require $2^{N(1-o(1))}$ distinct pattern sets, i.e. $N(1-o(1))$ bits
in the $nH=N/n$ denominator parameters.
\end{proof}


\begin{proposition}[Precision requirement]
\label{prop:wn-precision}
If every parameter of the normalized score \eqref{eq:wn-normal} is an
integer multiple of $2^{-p}$ of absolute value at most $2^p$, then at
most $2^{(2nH+1)(2p+2)}$ Boolean functions are representable with $H$
heads. Consequently a family of representations of all Boolean
functions with $H\le cN/n^2$ heads must use parameters with
$p\ge n/(4c)-1$ bits of precision.
\end{proposition}
\begin{proof}
Each parameter takes at most $2^{2p+1}+1\le2^{2p+2}$ values and the
function is determined by the parameters.
\end{proof}
With $p=O(\log n)$ (polynomial resolution) the bound is
$\Omega(N/(n\log n))$ heads and with constant precision $\Omega(N/n)$.
The spanning lists of Section~\ref{sec:upper-bounds} use integer
coefficients of size $O(n)$, consistent with their head counts.

\subsection{Rigidity of common directions}
\begin{theorem}[Common-direction collapse]
\label{thm:wn-common-direction}
Let $\beta\in\mathbb R^n_{\ge0}$, $s(x)=\beta\cdot x$, and suppose every
denominator has the form $B_h=1+t_h s$ with $t_h>0$ (or every denominator
has the form $1+t_h\,\beta\cdot(\mathbf1-x)$). Then every function
sign-represented by $\sigma+\sum_{h\le H}A_h/B_h$ with affine $A_h$ is a
sign pattern of the fixed linear space
\[
 \mathcal W=\Bigl\{R_0(s)+\sum_{i=1}^nx_iR_i(s):\ \deg R_0\le H,\ \deg R_i\le H-1\Bigr\},
 \qquad \dim\mathcal W\le(n+1)H+1 .
\]
Hence at most $2\sum_{i\le(n+1)H+1}\binom{N-1}i$ functions admit such a
representation, and representing every Boolean function (or a fraction
bounded away from zero of them) by lists with one common direction
requires $H\ge(\tfrac12-o(1))N/(n+1)$.
\end{theorem}
\begin{proof}
Multiply the score by $Q(s)=\prod_h(1+t_hs)>0$. Each
$A_h\prod_{g\ne h}(1+t_gs)$ equals $a_{h,0}\prod_{g\ne h}(1+t_gs)+\sum_ix_ia_{h,i}\prod_{g\ne h}(1+t_gs)$,
and the products have degree $\le H-1$ in $s$; $\sigma Q$ has degree
$H$. So $QS\in\mathcal W$, and $\operatorname{sign}S=\operatorname{sign}QS$.
The pattern count of a linear space of dimension $D$ on $N$ points is at
most $2\sum_{i\le D}\binom{N-1}i$, which is below $2^N$ unless
$D\ge(N-1)/2$ and is $o(2^N)$ unless $D\ge N/2-O(\sqrt N)$.
\end{proof}

The lists of Theorems~\ref{thm:uniform-spanning} and
\ref{thm:asymptotic-spanning} use $|D|\approx N/n$ distinct directions
$\mathbf1+v$, so they are not affected. The theorem shows that the
degree-$H$ factor in Warren's bound, which is the entire source of the
gap between $N/n^2$ and $N/n$, is illusory unless the denominator
directions spread: with one direction the products
$\prod_{g\ne h}B_g$ span a space independent of the $t_h$.

\subsection{A combinatorial sub-model and its transfer theorem}
\label{sec:wn-tropical}
Let head $h$ carry a permutation $\pi_h$ of $[n]$, distinct integers
$w_{h,1},\dots,w_{h,n}$ (distinct across all $nH$ pairs) and signs
$c_{h,k}\in\{\pm1\}$. For $x\ne0$ let $k_h(x)$ be the $\pi_h$-position of
the first element of $\operatorname{supp}(x)$ in the order $\pi_h$, and
put $T(x)=c_{h^\ast,k_{h^\ast}(x)}$ with
$h^\ast=\arg\max_hw_{h,k_h(x)}$. Equivalently $T$ is the decision list
whose rules are the $nH$ subcubes
$\Sigma_{h,k}=\{x_{\pi_h(k)}=1,\ x_{\pi_h(j)}=0\ (j<k)\}$, ordered by the
weights, each carrying a sign; the $n$ rules of one head partition
$\B^n\setminus\{0\}$.

\begin{theorem}[Tropical transfer]
\label{thm:wn-tropical}
If $T(x)=2f(x)-1$ for all $x\ne0$, then $\Hs(f)\le H$.
\end{theorem}
\begin{proof}
Put $L=2\max_{h,k}|w_{h,k}|+2$, $m_{h,\pi_h(k)}=kL$,
$e_{h,\pi_h(k)}=kL+w_{h,k}$, $b_{h,i}=K^{m_{h,i}}$ and
$u_{h,i}=c_{h,\pi_h^{-1}(i)}K^{e_{h,i}}$ for a large integer $K$. Fix
$x\ne0$, let $k=k_h(x)$ and $t=\pi_h(k)$. Every other element of the
support lies later in $\pi_h$, so its exponents are smaller by at least
$L/2$. Hence $u_h\cdot x=c_{h,k}K^{e_{h,t}}(1+\epsilon)$ and
$1+b_h\cdot x=K^{m_{h,t}}(1+\epsilon')$ with $|\epsilon|,|\epsilon'|\le nK^{-L/2}+K^{-L}$, and
\[
 \frac{u_h\cdot x}{1+b_h\cdot x}=c_{h,k}K^{w_{h,k}}(1+O(nK^{-L/2})).
\]
The weights are distinct integers, so the head attaining
$\max_hw_{h,k_h(x)}$ exceeds every other head by the factor $K$; for
$K>4nH$ the sign of $\sigma+\sum_hu_h\cdot x/(1+b_h\cdot x)$ equals
$T(x)$ for $x\ne0$ and equals $\operatorname{sign}\sigma$ at $x=0$. All
slopes are positive, so Theorem~\ref{thm:nf} applies.
\end{proof}

There are at most $(n!)^H(nH)!\,2^{nH}\le2^{nH(n+1)}$ tropical
representations with $H$ heads, so this sub-model cannot beat
$H=(1-o(1))N/n^2$; it is the regime of exponentially spread denominators
and it is a strict sub-model: parity needs about $N$ rules in it (each
fiber is a subcube minus earlier subcubes, hence an independent set of
the cube graph), whereas $\Hs(\mathrm{PARITY}_n)=n$. The cancellation
mechanism of Theorem~\ref{thm:sym} and the spanning mechanism of
Section~\ref{sec:upper-bounds} are therefore not captured by it.

\subsection{Rate caps of local mechanisms}
\begin{proposition}[Local cells cannot reach $N/n^2$]
\label{prop:wn-local-cap}
Suppose $k$ heads dominate a set $R\subseteq\B^n$ in the sense that on
$R$ their denominators are $O(s)$ and all other contributions are
$O(1)$, while off $R$ their contributions are bounded. Then as $s\to0$
the sign pattern on $R$ is a $k$-head pattern of the restricted problem,
hence one of at most $2(e|R|/(2nk))^{2nk}$ patterns. All $2^{|R|}$
patterns are available only if $|R|\le2nk\log_2(e|R|/(2nk))$. In
particular a mechanism that controls $|R|/k\ge n^2/c$ points per head
needs $|R|\ge2nk\,2^{n/(2c)}$: the dominated sets must have exponential
size.
\end{proposition}
\begin{proof}
On $R$ the score is $s^{-1}\sum_{j\le k}A_j/C_j+O(1)$ with $C_j$ the
positive affine forms $B_j/s$; for small $s$ its sign is that of
$\sum_jA_j/C_j$, a $k$-head score. Warren's theorem for $2nk$ parameters
and degree $k$ on $|R|$ points gives the count (Theorem~\ref{thm:worst}).
\end{proof}

For $k=1$ this is the star-cell mechanism of
Lemma~\ref{lem:pole-partition-spanning} ($n+1$ points per head); the
regular-zero mechanism of Theorem~\ref{thm:adaptive-cancellation} is
capped by its parameter count at $2n$ points per head, i.e. at
$W(n)\le(\tfrac12+o(1))N/n$, of which $2/3$ is attained. Boosting-type
arguments with the atoms as weak learners also stop at $N/n$, because
the atoms form a class of pseudo-dimension $O(n)$ and a random
function has correlation $O(n/\sqrt N)$ with every atom.


\subsection{The crux as an interpolation multiplicity}
\label{sec:wn-crux}
For fixed $\theta=(\sigma,u,b)$ the sign vector of $S(\cdot;\theta)$ is the
cell of $\theta$ in the arrangement of the $N$ hypersurfaces
$W_x=\{\theta:S(x;\theta)=0\}$ in $\mathbb R^{2nH}$. Warren's argument
bounds the number of cells by $\sum_{k\le2nH}\binom NkI_k$, where $I_k$
bounds the number of connected components of $k$-fold intersections; for
hyperplanes $I_k=1$ and the count is the linear one, while
Milnor--Thom gives $I_k\le(2\deg)^{2nH}$. A $(2nH)$-fold intersection is
the set of parameters whose score vanishes at $2nH$ prescribed cube
points, i.e.\ the solution set of a rational interpolation problem.
Hence $W(n)=\Theta(N/n^2)$ holds, up to constants in the exponent, if
and only if generic interpolation problems $S(x_j;\theta)=0$
($j\le2nH$) have $2^{\Theta(n^2H)}$ isolated real admissible solutions,
whereas $2^{O(nH\log n)}$ solutions would give $W(n)=\Omega(N/n)$. In
the univariate analogue the problem is linear after clearing
denominators (Cauchy interpolation) and has a unique solution; for
$H=1$ on the cube it is linear in $(u,b)$. For $H\ge2$ eliminating the
numerators leaves $nH$ polynomial conditions of degree up to $H(nH+1)$
in the $nH$ denominator parameters, whose B\'ezout number matches
Warren's bound; the number of real solutions in the admissible orthant
is the open quantity.


The univariate collapse has an algebraic explanation. The cleared score
is a first-order perturbation of the product $\prod_hB_h$, and products
of $H$ affine forms in $n$ variables form the Chow variety of completely
reducible forms, of dimension $nH$ and degree $(nH)!/(H!\,(n!)^H)$. For
$n=1$ every binary form splits, the Chow variety is the whole space, its
degree is one, and interpolation is unique. For $n\ge2$ Stirling's
formula gives $\log_2$ of the degree $\approx(n-1)H\log_2H$, which for
$H\asymp N/n^2$ is $\Theta(N)$: the complex multiplicity is
Warren-like. Hence no algebraic collapse occurs once the denominator
directions spread, and the open quantity is the number of real
admissible solutions for cube points.


\subsection{Computational evidence}
\label{sec:wn-experiments}
All searches below are heuristic: a success is certified by exact rational
evaluation of the score on every cube point, a failure proves nothing.
Scripts and logs are in \texttt{experiments/wn\_search/}.

\begin{itemize}
\item \emph{Alternating LP search.} For fixed denominators the best
numerators solve a linear program (maximum margin, or minimum hinge
loss); Danskin's theorem gives a subgradient in the denominators, which
are updated multiplicatively with random restarts. Random $7$-bit
functions are found with $8$ heads and random $8$-bit functions with
$12$ to $14$ heads, below the fixed-list threshold $nH\ge N/2$ of Cover's
bound ($16$ heads at $n=8$) but far above $N/n^2$. The certified
denominators have dynamic ranges $10^4$ to $10^{10}$ within a head and
spread directions. The same search does not recover the $n$-head
representation of parity for $n\ge6$, so its failures carry no
information.
\item \emph{Tropical sub-model.} An exact SAT decision procedure gives,
for random targets, $T(f)\in\{3,\dots,6\}$ at $n=5$ (mean $4.1$) and
$\{5,6,7\}$ at $n=6$ (mean $6.0$); parity needs $10$ heads at $n=5$ and
more than $13$ at $n=6$, majority more than $7$ at $n=6$. The optimal
small solutions catch one colour class exactly and let a default rule
take the rest, a structure that does not scale. A greedy constructor
with random orders needs $Hn/N\approx1.6,2.1,3.1,3.3$ heads at
$n=5,6,7,8$, i.e.\ more than $N/n$ and increasing.
\item \emph{Interpolation multiplicity.} For $n=5$, $H=2$ and a planted
target at $2nH=20$ cube points, Levenberg--Marquardt from $1500$ random
starts finds $19$ distinct real solutions, $3$ of them admissible: the
interpolation problem is not unique, in line with the Chow-degree
heuristic, but the sample is too small to measure growth.
\item \emph{Weight-two miniature.} Restricting the family to the
$\binom n2$ inputs of weight two (graphs), where Cover's threshold is
$H\approx n/4$ and Warren's about $n/(2\log_2n)$, the search finds
random graphs on $16$, $32$ vertices with $3$, $4$ heads; at $64$
vertices $3$ heads fail and the run is ongoing.
\end{itemize}

\subsection{Status}
\label{sec:wn-status}
The lower bound $(\tfrac12+o(1))N/n^2$ is Warren's count applied to the
$2nH$ real parameters of \eqref{eq:wn-normal}; by
Proposition~\ref{prop:wn-information} it is attained only if the sign
pattern depends on the denominators in a Warren-tight way. Every
mechanism available in the paper (exact spanning by fixed lists,
regular zeros, dominant cells, boosting) is capped at order $N/n$ by
the arguments above, and the tropical regime with random orders is
worse than $N/n$ in the greedy experiments. The analogous question for
depth-two threshold circuits (threshold of thresholds) is open with a
gap between the counting bound $\Omega(2^n/n^2)$ and the trivial
$2^{n-1}+1$; the only Warren-tight dictionary result on the cube known to
us is Selezneva's $\Theta(2^n/n^2)$ for exclusive-or sums of products of
affine forms over $\mathbb F_2$, whose mechanism (exact cancellation in
characteristic two) has no real-signed counterpart in the present model.
We state the two possibilities as a conjecture pair: either
$W(n)=\Theta(N/n^2)$, which requires a global construction in which
$\Theta(n)$ heads contribute comparably at a typical input, or
$W(n)=\omega(N/n^2)$, which requires a lower bound beyond Warren's
theorem that exploits the positivity and one-sign structure of the
denominators. Neither is established here.
